\uuid{DQYn}
\exo7id{4702}
\titre{exo7 4702}
\auteur{quercia}
\organisation{exo7}
\datecreate{2010-03-16}
\isIndication{false}
\isCorrection{true}
\chapitre{Suite}
\sousChapitre{Convergence}
\cpge{mpsi-mp2i.nombres-reels-suites.c, pcsi.nombres-reels-suites.b, ptsi.nombres-reels-suites.b, mpsi-mp2i.nombres-reels-suites.e, pcsi.nombres-reels-suites.d, ptsi.nombres-reels-suites.d, mpsi-mp2i.nombres-reels-suites.d, pcsi.nombres-reels-suites.c, ptsi.nombres-reels-suites.c, bcpst1.suites-reelles, bcpst1.limites-continuite.b}
\module{Analyse}
\niveau{L1}
\difficulte{}

\contenu{
\question{
Soit $u_0\in\N^*$ et $(u_n)$ la suite d{\'e}finie par la relation
de r{\'e}currence~: $u_{n+1} = u_n^2+1$. Montrer qu'il exitste $a\in\R$ tel
que $u_n = [a^{2^n}]$ pour tout~$n$ o{\`u} $[\ ]$ d{\'e}signe la partie enti{\`e}re.
}
\reponse{
Les suites $(\ln(u_n)/2^n)$ et $(\ln(1+u_n)/2^n)$ sont adjacentes.
}
}
